\documentclass[10pt,a4paper]{amsart}
\usepackage[margin=19mm]{geometry}
\usepackage{amsmath,amssymb,mathtools}
\usepackage[scr=rsfs,cal=euler]{mathalfa}
\usepackage{xfrac}
\usepackage[unicode,hidelinks,bookmarks=false]{hyperref}
\newcommand{\ie}{i.e.\ }
\newcommand{\cf}{cf.\ }
\newcommand{\resp}{respectively\ }
\newcommand{\Subsection}{\subsection}
\providecommand{\nobreakdash}{\nobreak\hbox{-}}
\setlength{\emergencystretch}{2em}
\multlinegap0pt
\usepackage{bm}
\usepackage{bbm}
\usepackage{dsfont}
\usepackage{xspace}
\usepackage{cancel}
\usepackage{mathtools}
\usepackage{amsthm}
\makeatletter
\@ifundefined{lemme}{\newtheorem{lemme}{Lemma}}{}
\@ifundefined{prop}{\newtheorem{prop}{Proposition}}{}
\makeatother
\newcommand{\Fke}{F_{k,\ve}}
\newcommand{\Om}{\Omega}
\newcommand{\R}{\mathbb{R}}
\newcommand{\Sp}{{\rm Sp}}
\newcommand{\mkle}{\mu_{k,\ve}}
\newcommand{\mkl}{\mu_{k}^{\lambda}}
\newcommand{\mpl}{\mu_{p}^{\lambda}}
\newcommand{\ve}{\varepsilon}
\newcommand{\vil}{\varphi_{i}^{\lambda}}
\newcommand{\vjl}{\varphi_{j}^{\lambda}}
\newcommand{\vkl}{\varphi_{k}^{\lambda}}
\newcommand{\vp}{\varphi}
\title[Least-energy solutions of the Br\'ezis-Nirenberg problem in ]{Least-energy solutions of the Br\'ezis-Nirenberg problem in the non-coercive case in dimension $3$}
\author{Hussein Cheikh Ali, Bruno Premoselli}
\date{Source: arXiv:2509.19145v1 (CC BY 4.0)}
\begin{document}
\maketitle

\noindent\emph{Source and licence.} Least-energy solutions of the Br\'ezis-Nirenberg problem in the non-coercive case in dimension $3$. Version arXiv:2509.19145v1 (CC BY 4.0), distributed under the Creative Commons Attribution 4.0 International licence, as recorded in the arXiv licence metadata for this exact version and shown on the abstract page https://arxiv.org/abs/2509.19145v1. Full rights notice: licenses/arxiv-2509.19145-CC-BY-4.0.txt.\par\smallskip

\noindent\textit{Editorial scope.} Counted unit: the Lemme whose source label is \texttt{lemme2} in the article (source0.tex lines 926--928, source label \texttt{lemme2}), delivered together with the article's own complete original proof of it (source0.tex lines 930--990, one \texttt{proof} environment of 61 lines). No mathematics is written, repaired or re-derived: statement and proof are the article's own text line for line. Required context retained uncounted: defK0 (lines 182--184); muk (lines 374--376); prop:vp (lines 456--484); vpk (lines 460--462); prop:signvp (lines 554--563); holder:spectre (lines 589--591); mukeps (lines 897--899); eq:test1e (lines 913--915). Every definition, display and label of those lines is retained unchanged. Omitted from this package and not redistributed: 1--181, 185--373, 377--455, 463--553, 564--588, 592--896, 900--912, 916--925, 929--929, 991--2096. Nothing inside a retained excerpt is omitted or altered: every retained body is delivered exactly as the article's lines are written, so no omission receipt of any kind accompanies this package. Citations: the retained excerpts carry no citation command at all, so no bibliography entry of the article is transcribed and no reference is redistributed. Reference adaptation: none was needed and none was made; every reference inside the delivered unit resolves inside it, so no unresolved reference or citation is produced. Printed slips kept literal: no printed word, symbol, hypothesis or number of the counted statement or of its proof was altered. Layout and class: the article's own class is \texttt{amsart} with packages including amssymb, graphicx, dsfont, xcolor, hyperref; the wrapper is the lane's pinned \texttt{amsart} editorial wrapper at 10pt on A4, which restates the theorem-like environments the retained text needs and adds the article's own macro definitions that the retained text needs (\texttt{\textbackslash Fke}, \texttt{\textbackslash Om}, \texttt{\textbackslash R}, \texttt{\textbackslash Sp}, \texttt{\textbackslash mkl}, \texttt{\textbackslash mkle}, \texttt{\textbackslash mpl}, \texttt{\textbackslash ve}, \texttt{\textbackslash vil}, \texttt{\textbackslash vjl}, \texttt{\textbackslash vkl}, \texttt{\textbackslash vp}), with extra wrapper packages (bm, bbm, dsfont, xspace, cancel, mathtools). The delivered numbering is the unit's own. Assets: no retained span contains a figure environment, a graphics-inclusion command, a PDF-page inclusion, a table or a TikZ picture; no image file is copied into this package, the package declares no asset dependency and no image file is redistributed with it. Licence: version arXiv:2509.19145v1 of this article is distributed under the Creative Commons Attribution 4.0 International licence, as recorded in the arXiv licence metadata for this exact version and shown on the abstract page https://arxiv.org/abs/2509.19145v1. A full rights notice is delivered at licenses/arxiv-2509.19145-CC-BY-4.0.txt. Characters: every retained excerpt is pure ASCII, so the delivered engine, XeLaTeX with its default Latin Modern text font, reports no missing character.
\begin{equation} \label{defK0}
 \frac{1}{K_n^2} = \inf_{\vp \in C^\infty_c(\R^n) \backslash \{0\}} \frac{\int_{\R^n}|\nabla \vp|^2 dx}{\big(\int_{\R^n} |\vp|^{\frac{2n}{n-2}}dx \big)^{\frac{n-2}{n}}}. 
 \end{equation}

\begin{equation} \label{muk}
\mkl(u) = \inf_{\text{dim}_u V = k} \max_{v \in V \backslash \{0\}} \frac{\int_{\Om}\left(  |\nabla v|^2 - \lambda v^2 \right) \, dx}{\int_{\Om} u^{2^*-2} v^2 \, dx},
\end{equation}

\begin{prop} \label{prop:vp}
Let $\Om \subset \R^n$, $n \ge 3$, be bounded and smooth. Let $\lambda \ge0$ and $u \in L^{2^*}_>(\Om)$. Then
\begin{enumerate}
\item For any $k \ge 1$, $\mkl(u)> - \infty$ and it is attained at some generalised eigenfunction $\vkl(u) \in H^1_0(\Om) \backslash \{0\}$ that satisfies weakly
\begin{equation} \label{vpk}
\big( - \Delta - \lambda \big) \vkl(u) = \mu_{k}^{\lambda}(u) u^{2^*-2}    \vkl(u) \quad \text{ in } \Om.
\end{equation}
The functions $(\vkl(u))_{k \ge 1}$ are normalised to satisfy
\begin{equation} \label{norma} 
\int_\Om u^{2^*-2} \vil(u) \vjl(u) dx =  \delta_{ij} \quad \text{ for } i \neq j,
\end{equation}
and for $k \ge 2$ each $\vkl(u)$ is sign-changing. 
\item The following equivalent characterisation of $\mkl(u)$ holds: for any $k \ge 1$,
\begin{equation} \label{muk2}
\mkl(u) = \inf_{v \in S_{k-1}} \frac{\int_{\Om}\left(  |\nabla v|^2 - \lambda v^2 \right) \, dx}{\int_{\Om} u^{2^*-2} v^2 \, dx} 
\end{equation}
where 
$$\begin{aligned} 
S_{k-1} & = \Big\{ v \in H^1_0(\Om)  \backslash \{0\} \text{ such that } \int_\Om u^{2^*-2} \vil(u) v dx = 0 \\
&\text{ for all } 1 \le i \le k-1 \Big\}.
\end{aligned} $$
\item For $k \ge 1$ we denote by $E_k^{\lambda}(u) \subset H^1_0(\Om)$ the vector space of solutions of \eqref{vpk}. Then $\dim E_k^{\lambda}(u) < + \infty$.

\item We have $\mkl(u) \le \mu_{k+1}^{\lambda}(u)$ for any $k \ge 1$ and there exists $\gamma = \gamma (u) >0$ such that $ \mu_1^{\lambda}(u) + \gamma \le \mu_{2}^{\lambda}(u)$.
\item We have $\mkl(u) \to + \infty$ as $k \to + \infty$. 
\item If there exist $\vp \in H^1_0(\Om) \backslash \{0\}$ and $\mu \in \R$ such that $\big( - \triangle - \lambda \big) \vp = \mu u^{2^*-2} \vp$ in  $\Om$ then $\mu =\mkl(u)$ for some $k \ge 1$ and $\vp \in E_k^{\lambda}(u)$. 

 \end{enumerate}
\end{prop}

 \begin{prop}  \label{prop:signvp}
  Let $\Om \subset \R^n$, $n \ge 3$, be bounded and smooth and let  $\lambda \ge 0$. Let $i \ge0$ be such that $\lambda \in [\Lambda_i, \Lambda_{i+1})$ and let $k = N(i)+1 \ge 1$. We have  
\begin{equation} \label{infmukpos}
 \inf_{u \in L^{2^*}_{>}(\Om)} \mu_k^{\lambda}(u)  \Vert u \Vert_{2^*}^{2^*-2}  > 0
 \end{equation}
and, for any $1 \le p \le k-1$,
\begin{equation} \label{supmupneg}
\sup_{u \in L^{2^*}_{>}(\Om)} \mpl(u)  \Vert u \Vert_{2^*}^{2^*-2} \le 0 .
\end{equation}
\end{prop}

\begin{equation} \label{holder:spectre}
 \Vert u_m^{2^*-2} \vp_{k,m} \Vert_{\frac{2n}{n+2}}^{\frac{2n}{n+2}} \le \Big(\int_{\Om} u_m^{2^*}dx\Big)^{\frac{2}{n+2}} \Big( \int_{\Om} u_m^{2^*-2} \vp_{k,m}^2 dx\Big)^{\frac{n}{n+2}} = 1, 
 \end{equation}

\begin{equation} \label{mukeps}
\mkle = \inf_{u \in X_\ve, \Vert u \Vert_{2^*}=1} \Fke(u). 
\end{equation}

\begin{equation} \label{eq:test1e}
\mkle < K_n^{-2}.
\end{equation}

\begin{lemme} \label{lemme2}
Let $\ve_0 >0$ be such that, for every $0 < \ve \le \ve_0$, \eqref{eq:test1e} is satisfied. Then, for all $0 < \ve \le \ve_0$, $\mkle$ is attained at some function $u_\ve \in X_\ve $ with $\Vert u_\ve \Vert_{2^*} = 1$.   
\end{lemme}

\begin{proof}
We assume that $\ve>0$ is fixed small enough so that \eqref{eq:test1e} is satisfied. We let $(u_m)_{m\ge 1}$ be a minimising sequence for $\mkle$: $u_m \in X_\ve$, $\Vert u_m \Vert_{2^*} =1$ and 
\begin{equation} \label{eq:cve1}
\Fke(u_m) = \mkl(u_m)  + \ve \int_{\Om} u_m^{-\ve}dx  \to \mkle
\end{equation}
 as $m \to + \infty$. For every $m \ge 1$, $u_m \in L^{2^*}_>(\Om)$, so by Proposition \ref{prop:vp} a generalised eigenfunction $\vkl(u_m)$ exists. We will let for simplicity in what follows $\vp_{k,m} = \vkl(u_m)$ and $ \mu_{k,m} = \mkl(u_m)$. This function satisfies $\int_{\Om} u_m^{2^*-2} \vp_{k,m}^2dx = 1$  and 
\begin{equation} \label{eq:cve2}
 - \Delta \vp_{k,m} - \lambda \vp_{k,m} = \mu_{k,m} u_m^{2^*-2} \vp_{k,m} \quad \text{ in } \Om. 
 \end{equation}
Let in what follows $w_m = u_m^{2^*-2}$. Then $(w_m)_{m \ge1}$ is uniformly bounded in $L^{\frac{n}{2}}(\Om)$ and there exists $w_\ve \in L^{\frac{n}{2}}(\Om)$, $w_\ve \ge 0$ a.e., such that $w_m \rightharpoonup w_\ve$ weakly in $L^{\frac{n}{2}}(\Om)$, up to a subsequence, as $m \to + \infty$. Let $u_\ve = w_\ve^{\frac{1}{2^*-2}} \ge0$ a.e. By lower semi-continuity of the $L^{\frac{n}{2}}(\Om)$ norm we have $ \int_{\Om} u_\ve^{2^*}dx \le 1$. Since $t \mapsto t^{-\frac{\ve}{2^*-2}}$ is convex in $(0, + \infty)$, $u \mapsto \int_{\Om} u^{-\frac{\ve}{2^*-2}} dx$ is weakly lower semi-continuous over $\{ w \in L^{\frac{n}{2}}(\Om), w \ge 0 \text{ a.e.} \}$. As a consequence of the weak convergence and of \eqref{eq:cve1} we get
\begin{equation} \label{eq:cve3}
\begin{aligned}
 \int_{\Om} u_\ve^{-\ve}dx = \int_{\Om} {w_\ve}^{- \frac{\ve}{2^*-2}} &\le \liminf_{m \to + \infty}  \int_{\Om} w_m^{-\frac{\ve}{2^*-2}}dx \\
 &=  \liminf_{m \to + \infty}  \int_{\Om} u_m^{-\ve}dx  \le \frac{1}{\ve} \mkle < + \infty. 
\end{aligned} 
 \end{equation}
This shows that $u_\ve \in X_\ve$ and hence that $u_\ve \neq 0$. By \eqref{eq:cve1} the sequence $(\mu_{k,m})_{m \ge1}$ is bounded. We claim that there exists $C >0$ such that 
\begin{equation} \label{eq:bornevpkm}
\Vert \vp_{k,m} \Vert_{H^1_0} \le C
\end{equation}
for all $m \ge 1$. First, if $\lambda \in (\Lambda_i, \Lambda_{i+1})$, that is if $\lambda \not \in \Sp(-\Delta)$, $-\Delta- \lambda$ is invertible and \eqref{eq:bornevpkm} follows from \eqref{eq:cve2} since the right-hand side of \eqref{eq:cve2} is uniformly bounded in $L^{\frac{2n}{n+2}}(\Om)$ (this is the same argument that we used in \eqref{holder:spectre}). We may thus assume that $\lambda = \Lambda_i$. We proceed as in the proof of Proposition \ref{prop:signvp}: We split $\vp_{k,m}$ as 
$$ \vp_{k,m} = v_m + z_m, $$
where $z_m \in E_{\Lambda_i}(\Om)$ and $v_m$ is $H^1_0(\Om)$-orthogonal to $E_{\Lambda_i}(\Om)$. Since $-\Delta z_m - \Lambda_i z_m = 0$, by \eqref{eq:cve2} $v_m$ satisfies 
\begin{equation*} 
 - \Delta v_m - \Lambda_i v_m = \mu_{k,m} u_m^{2^*-2} \vp_{k,m} \quad \text{ in } \Om. 
 \end{equation*}
By definition $- \Delta-\Lambda_i$ is invertible in $E_{\Lambda_i}(\Om)^{\perp}$ and the same argument as above shows that $\Vert v_m \Vert_{H^1_0} \le C $ for some $C>0$ and for all $m \ge 0$. To prove \eqref{eq:bornevpkm} we thus proceed by contradiction and assume that $\Vert z_m \Vert_{H^1_0} \to + \infty$ as $m \to + \infty$, and we let $\tilde{\vp}_{k,m} = \frac{\vp_{k,m}}{\Vert z_m\Vert_{H^1_0}}$. Since $(v_m)_{m \ge0}$ is bounded in $H^1_0(\Om)$ we have 
$$ \tilde{\vp}_{k,m} = \frac{z_m}{\Vert z_m\Vert_{H^1_0}} + \frac{v_m}{\Vert z_m\Vert_{H^1_0}} = \frac{z_m}{\Vert z_m\Vert_{H^1_0}}  + o(1) \quad \text{ in } H^1_0(\Om). $$
The function  $ \tilde{z}_m:=\frac{z_m}{\Vert z_m\Vert_{H^1_0}}$ is bounded in $H^1_0(\Om)$ and thus weakly converges to some $Z \in H^1_0(\Om)$. Since $\tilde{z}_m \in E_{\Lambda_i}(\Om)$ for every $m \ge 0$ and since $E_{\Lambda_i}(\Om)$ is finite-dimensional, $\tilde{z}_m$ strongly converges to $Z$ in $H^1_0(\Om)$, and we have $\Vert Z \Vert_{H^1_0} =1$ and $Z \in E_{\Lambda_i}(\Om)$. Thus $\tilde{\vp}_{k,m}$ strongly converges to $Z$ in $H^1_0(\Om)$ as $m \to + \infty$. The normalisation condition $\int_{\Om} u_m^{2^*-2} \vp_{k,m}^2dx = 1$ now becomes:
\begin{equation} \label{eq:bornevpkm:2}
 \frac{1}{\Vert z_m \Vert_{H^1_0}^2} =  \int_{\Om} u_m^{2^*-2} \tilde{\vp}_{k,m}^2dx = \int_{\Om} u_m^{2^*-2} Z^2dx + o(1). 
 \end{equation}
Since $u_m^{2^*-2} \rightharpoonup u_{\ve}^{2^*-2}$ in $L^{\frac{n}{2}}(\Om)$ we can pass \eqref{eq:bornevpkm:2} to the limit  and we obtain 
$$ \int_{\Om} u_\ve^{2^*-2} Z^2dx = 0. $$
Since $u_\ve \in X_\ve$ we have in particular $u_\ve \in L^{2^*}_{>}(\Om)$, which implies that $Z = 0$ a.e. in $\Om$. But since $Z \in E_{\Lambda_i}(\Om)$ standard elliptic theory shows that it is smooth, and hence $Z = 0$, which contradict $\Vert Z \Vert_{H^1_0} = 1$. This proves \eqref{eq:bornevpkm}. 


We denote by $\vp_{k,\ve}$ the weak limit in $H^1_0(\Om)$, up to a subsequence, of $( \vp_{k,m} )_{m \ge 1}$, and we let $\hat{\mu}_{k,\ve} = \lim_{m \to +\infty} \mu_{k,m}$. Since $\vp_{k,m}$ strongly converges to $\vp_{k,\ve}$ in $L^\frac{n}{n-2}(\Om)$ by Sobolev's embedding and since $w_m$ weakly converges to $w_\ve = u_\ve^{2^*-2}$ in $L^{\frac{n}{2}}(\Om)$ we get that $\vp_{k,\ve}$ weakly satisfies
 \begin{equation} \label{eq:cve4}
  - \Delta \vp_{k,\ve} - \lambda \vp_{k,\ve} = \hat{\mu}_{k,\ve}  u_\ve^{2^*-2}\vp_{k,\ve} \quad \text{ in } \Om. 
 \end{equation}
We claim that $\vp_{k,\ve} \neq 0$. Assume otherwise: then $\vp_{k,m} \to 0$ in $L^2(\Om)$ up to a subsequence as $m \to + \infty$. Integrating \eqref{eq:cve2} against $\vp_{k,m}$ and using H\"older and Sobolev inequalities yields, since $\Vert u_m \Vert_{2^*} = 1$,
$$ \begin{aligned} \int_{\Om} |\nabla \vp_{k,m}|^2dx + o(1) & \le \mu_{k,m} \big( \int_{\Om} |\vp_{k,m}|^{2^*}dx \big)^{\frac{n-2}{n}}  \le \mu_{k,m}K_n^2  \int_{\Om} |\nabla \vp_{k,m}|^2dx, \end{aligned} $$
where $K_n$ is the optimal constant in Sobolev's inequality given by \eqref{defK0}. By \eqref{eq:test1e} we have $\limsup_{m \to  + \infty} \mu_{k,m} < K_n^{-2}$, so the latter yields $\int_{\Om} |\nabla \vp_{k,m}|^2 dx = o(1)$. But this is a contradiction with $\int_{\Om} u_m^{2^*-2} \vp_{k,m}^2dx = 1$, hence $\vp_{k,\ve} \neq 0$. 


We now remark that $\hat{\mu}_{k,\ve} >0$. Indeed, we have $\hat{\mu}_{k,\ve} =  \lim_{m \to + \infty} \mu_{k,m}$ (up to a subsequence), and $\mu_{k,m} = \mkl(u_m) \ge \mkl(\Om) >0$, where we used Proposition \ref{prop:signvp}. By Proposition \ref{prop:signvp}, and since $\lambda \in [\Lambda_i, \Lambda_{i+1})$, the $N(i)$ first generalised eigenvalues $\mu_1^{\lambda}(u_\ve), \dots, \mu_{N(i)}^{\lambda}(u_\ve)$ are nonpositive. Since $\hat{\mu}_{k,\ve} >0$, \eqref{eq:cve4} implies that 
$$ \int_{\Om} u_\ve^{2^*-2} \vp_{k,\ve} \vp_p^{\lambda}(u_\ve)dx = 0 \quad \text{ for all } 1 \le p \le k-1. $$
Thus $V =  \text{Span} \big( \vp_1^{\lambda}(u_\ve), \dots, \vp_{k-1}^\lambda(u_\ve), \vp_{k,\ve}\big) $ satisfies $\dim_{u_\ve} V = k$ and is an admissible subspace in the definition of $\mkl(u_\ve)$: \eqref{muk} and \eqref{eq:cve4} then show that 
\begin{equation} \label{eq:cve5}
\mkl(u_\ve) \le \hat{\mu}_{k,\ve}.
\end{equation}
 It remains to let $ \tilde{u}_\ve = \frac{u_\ve}{\Vert u_\ve \Vert_{2^*}}$, so that $\Vert \tilde{u}_\ve \Vert_{2^*} = 1$. Since $\Vert u_\ve \Vert_{2^*} \le 1$ we have, using \eqref{eq:cve1},  \eqref{eq:cve3}, \eqref{eq:cve5} for $\ve$ small enough:
$$ \begin{aligned}
\Fke(\tilde{u}_\ve) & =  \mkl(\tilde{u}_\ve) + \ve  \Big(\int_{\Om} \tilde{u}_\ve^{-\ve}dx \Big)\\
& = \mkl(u_\ve) \Vert u_\ve \Vert_{2^*}^{2^*-2} + \ve \Vert u_\ve \Vert_{2^*}^{\ve} \int_{\Om}u_\ve^{-\ve}dx \\
& \le \Vert u_\ve \Vert_{2^*}^{\ve}\Big( \hat{\mu}_{k,\ve} + \ve \int_{\Om} u_\ve^{-\ve}dx \Big) \\
& \le \Vert u_\ve \Vert_{2^*}^{\ve}\liminf_{m \to + \infty} \Big(  \mkl(u_m) + \ve \int_{\Om} u_m^{-\ve}dx \Big)  = \Vert u_\ve \Vert_{2^*}^{\ve} \mkle.
\end{aligned} $$
By \eqref{mukeps} this shows at once that $\Vert u_\ve \Vert_{2^*} = 1$ and that $u_\ve$ attains $\mkle$. 
\end{proof}

\end{document}
